%
% raender.tex -- Definition für Raender
%
% (c) 2021 Prof Dr Andreas Müller, OST Ostschweizer Fachhochschule
%
\documentclass[tikz]{standalone}
\usepackage{amsmath}
\usepackage{times}
\usepackage{txfonts}
\usepackage{pgfplots}
\usepackage{csvsimple}
\usetikzlibrary{arrows,intersections,math,calc}
\begin{document}
\def\skala{1}
\definecolor{darkred}{rgb}{0.8,0,0}
\def\punkt#1#2{
	\fill[color=white] #1 circle[radius=0.05];
	\draw[color=#2,line width=1.2pt] #1 circle[radius=0.05];
}
\begin{tikzpicture}[>=latex,thick,scale=\skala]

\coordinate (A) at (1,2);
\coordinate (B) at (4,1);
\coordinate (C) at (5,4);
\coordinate (D) at (2,5);

\coordinate (A0) at (-6,1);
\coordinate (B0) at (-2,1);
\coordinate (C0) at (-2,5);
\coordinate (D0) at (-6,5);

\fill[color=gray!20] (A0) rectangle (C0);
%\draw[line width=1.2pt] (A0) rectangle (C0);
\draw[->,line width=1.2pt] (A0) -- (B0);
\draw[->,line width=1.2pt] (B0) -- (C0);
\draw[->,line width=1.2pt] (C0) -- (D0);
\draw[->,line width=1.2pt] (D0) -- (A0);
\node at ($0.5*(A0)+0.5*(C0)$) {$I\times I$};

\node at (A0) [below] {$(0,0)$};
\node at (B0) [below] {$(1,0)$};
\node at (C0) [above] {$(1,1)$};
\node at (D0) [above] {$(0,1)$};

\punkt{(A0)}{black}
\punkt{(B0)}{black}
\punkt{(C0)}{black}
\punkt{(D0)}{black}

\draw[->,color=blue,shorten <= 0.2cm,shorten >= 0.2cm]
	(A0) to[out=30,in=160] (A);
\draw[->,color=blue,shorten <= 0.2cm,shorten >= 0.2cm]
	(B0) to[out=-10,in=-150] (B);
\draw[->,color=blue,shorten <= 0.2cm,shorten >= 0.2cm]
	(C0) to[out=20,in=130] (C);
\draw[->,color=blue,shorten <= 0.2cm,shorten >= 0.2cm]
	(D0) to[out=-20,in=-160] (D);

\node at (A) [below left] {$A$};
\node at (B) [below right] {$B$};
\node at (C) [above right] {$C$};
\node at (D) [above left] {$D$};

\fill[color=darkred!20] (A)
	to[out=-10,in=-170] (B)
	to[out=90,in=-90] (C)
	to[out=-180,in=10] (D)
	to[out=-100,in=100] cycle;
\draw[->,color=darkred,line width=1.2pt] (A) to[out=-10,in=-172] (B);
\draw[->,color=darkred,line width=1.2pt] (B) to[out=90,in=-91] (C);
\draw[->,color=darkred,line width=1.2pt] (C) to[out=-180,in=7] (D);
\draw[->,color=darkred,line width=1.2pt] (D) to[out=-100,in=98] (A);
%\draw[color=darkred,line width=1.2pt] (A)
%	to[out=-10,in=-170] (B)
%	to[out=90,in=-90] (C)
%	to[out=-180,in=10] (D)
%	to[out=-100,in=100] cycle;

\node at ($0.25*(A)+0.25*(B)+0.25*(C)+0.25*(D)$) {$\Gamma(I\times I)$};

\punkt{(A)}{darkred}
\punkt{(B)}{darkred}
\punkt{(C)}{darkred}
\punkt{(D)}{darkred}

\node[color=darkred] at ($0.5*(A)+0.5*(B)$) [below left]  {$\gamma_1$};
\node[color=darkred] at ($0.5*(B)+0.5*(C)$) [below right] {$\gamma_2$};
\node[color=darkred] at ($0.5*(C)+0.5*(D)$) [above right] {$\gamma_3$};
\node[color=darkred] at ($0.5*(D)+0.5*(A)$) [above left]  {$\gamma_4$};

\draw[->] (-0.1,0) -- (5.8,0)
	coordinate[label={$\operatorname{Re}$}];
\draw[->] (0,-0.1) -- (0,6.3)
	coordinate[label={right:$\operatorname{Im}$}];

\node at (5.5,6) {$\mathbb{C}$};

\end{tikzpicture}
\end{document}

